\documentclass[11pt]{article}
\usepackage[margin=1in]{geometry}
\usepackage{times}
\usepackage{hyperref}
\hypersetup{colorlinks=true, linkcolor=black, urlcolor=blue, citecolor=black}
\title{The Ball That Fell Faster: Bruce DePalma's Spin Experiments and the Timing Problem}
\author{Flyxion \\ \small Independent Researcher}
\date{}
\begin{document}
\maketitle

\begin{abstract}
Bruce DePalma's reported 1970s experiments, in which a rapidly spinning steel ball launched vertically was claimed to rise higher and fall in less time than an identical non-spinning ball launched with the same initial velocity, have circulated as evidence that rotation can alter an object's inertial or gravitational mass. This paper argues that the reported result is best explained by systematic timing and launch-mechanism asymmetries between the spinning and non-spinning trials rather than by any genuine alteration of gravitational behavior, that the closest controlled tests of the underlying idea, in which spinning rotors were weighed directly, found no dependence of weight on rotation, and that the proposed mechanism, insofar as DePalma specified one, does not identify any established channel by which mechanical spin angular momentum could couple to gravitational mass. The paper also places the spin result beside DePalma's better known homopolar generator claims, since both rest on the same style of argument. It does not offer a new measurement of a spinning ball.
\end{abstract}

\section{Introduction}

Bruce DePalma, an electrical engineer trained at MIT, reported in the mid-1970s that a steel ball spun to high angular velocity and launched vertically from a spring-loaded mechanism rose noticeably higher, and returned to the launch point in noticeably less time, than an otherwise identical non-spinning ball launched from the same mechanism at nominally the same speed. DePalma interpreted this as evidence that rapid spin reduces an object's effective gravitational mass, a claim he connected to a broader proposed theory of rotational physics distinct from standard mechanics. His 1974 monograph on the generation of a unidirectional force sets out the program, and his later papers extend it to the extraction of electrical energy from space by means of a rotating magnetized conductor, the device he called the N machine \cite{DP1974,DP1990}.

The spinning ball has a place in the antigravity literature because it is simple. It needs no vacuum chamber and no high voltage, and a timing result of the kind reported can be described in a sentence. It is that simplicity that makes the experiment a useful test case. A simple experiment makes few demands on the apparatus, and so the failure of the control that was not run, the measurement of the launch speed, is easy to identify.

\section{The Launch Mechanism Problem}

The reported experiment compared a spinning ball to a non-spinning ball launched from the same spring mechanism, but a ball's interaction with a mechanical launcher is not spin-independent. A rapidly spinning ball in contact with a launch cup or track experiences frictional and gyroscopic interactions during the launch phase itself that a non-spinning ball does not, potentially imparting a different effective launch velocity even when the spring mechanism's stored energy is nominally identical between trials. Spin couples to the surface it touches, and a ball released from contact while rotating can leave with a different translational speed than a ball released from the same contact without rotation, either from a transfer of angular to linear momentum at the surface or from a difference in how the spring's force is delivered to a rolling body.

DePalma's original reports do not, as far as the secondary literature records, describe a launch-velocity measurement independent of the timing results themselves, meaning the comparison assumes, rather than independently confirms, that both balls left the launcher at the same speed. The assumption is exactly the one a spin-dependent launch interaction would violate. A difference in flight time of the size reported corresponds to a difference in launch speed of a few percent, well inside the range that a mechanical launcher can produce from run to run. Establishing a gravitational effect would require showing that the difference in flight time persists when the difference in launch speed is measured and removed.

\section{What Controlled Testing of the Underlying Idea Has Found}

The claim that spin changes weight was tested directly on more than one occasion. In 1989 Hayasaka and Takeuchi reported that a spinning gyroscope appeared to lose weight when it was rotated in one sense, a result that drew wide attention \cite{HT1989}. The following year two groups reported independent tests with improved instruments. Faller and colleagues weighed a spinning gyroscope and found a null result \cite{F1990}, and Quinn and Picard measured the mass of spinning rotors and found no dependence on rotation speed or sense of rotation \cite{QP1990}. Their conclusion was that the effect did not exist at the level of the earlier report. These are the closest controlled tests of the central idea in the DePalma claim, that rotation alters weight, and the outcome is a null.

At a higher level of precision, the equivalence of inertial and gravitational mass has been tested to a few parts in $10^{15}$ in the final results of the MICROSCOPE satellite experiment \cite{T2022}. Those tests compare test masses of different composition, not spinning bodies, so they bear on the general principle and not on the particular claim. They do show that the framework in which a spinning ball would behave differently from a stationary one has to survive an extraordinarily tight constraint on the underlying coupling of gravity to matter.

I could not locate a published replication of the spinning ball launch itself that used independent measurement of launch velocity, such as high-speed photography of the balls leaving the launcher. The absence matters in both directions. It means the strongest statement available is not that the specific launch experiment has been refuted by a specific replication, but that the original comparison lacked the control that would make it informative, and that tests of the physical claim in forms that could be controlled, the weighed rotors, returned nulls.

\section{Why the Proposed Mechanism Was Never Specified Precisely Enough to Test Independently}

DePalma's broader theoretical writing proposed that rotation constitutes a distinct and previously unrecognized form of energy and inertia not captured by standard mechanics, but did not derive a quantitative prediction for how much gravitational mass reduction a given angular velocity should produce, a prediction that would be needed to distinguish the claim from a null result on statistical grounds rather than relying on a single, methodologically compromised timing comparison. Without a quantitative prediction, the claim cannot be said to have survived null replications so much as to have never been placed in a position to be tested against them precisely.

The same pattern appears in the N machine. Kincheloe examined the Sunburst N machine in the mid-1980s and, in his report, observed that DePalma's 1980 power measurements gave no consideration to the phase angle of the alternating-current drive as the motor load increased, which gives optimistic results, since power in an a.c. circuit cannot be computed without it \cite{K1986}. A discrepancy of this kind is what a quantitative theory with stated measurement requirements would let one anticipate. Without the theory, each result is merely a number to be explained after the fact.

\section{The Entropic Gravity Framing}

On an account where gravitational curvature tracks the redistribution of an ordered configuration at scales where mass-energy content is significant, the angular momentum of a spinning steel ball a few centimeters across represents a change in the ball's internal organization utterly negligible relative to what would be required to alter its participation in that redistribution, consistent with, though not dependent on, the general-relativistic expectation that inertial and gravitational mass are equivalent regardless of an object's internal rotational state, a form of the equivalence principle tested to high precision by entirely different experiments with no observed spin dependence \cite{T2022}.

\section{Objections}

It may be said that the weighed-rotor experiments tested a different quantity from the launched ball, since weight in a balance is not the same as acceleration in free flight. This is true, and it is why the tests are described here as the closest controlled tests and not as direct replications. But the physical claim in both cases is the same, that spin changes the gravitational response of a body, and a mechanism that altered free-fall acceleration without altering weight would be a violation of the equivalence principle of a stronger kind than the claim requires.

It may also be said that the timing difference was seen repeatedly by the experimenter, and that repetition is evidence. Repetition of an uncontrolled comparison repeats the control failure. A systematic difference in launch speed would produce the same difference in flight time on every trial.

\section{Conclusion}

The DePalma spinning ball result is a case study in how an uncontrolled variable in a launch mechanism can produce a striking apparent effect, and the absence of a quantitative theoretical prediction from DePalma's own proposed mechanism means the claim was never precise enough to be meaningfully confirmed or refuted by a single demonstration in the first place. The better controlled tests of whether spin changes weight, in the weighed-rotor experiments of 1990, found no effect. This paper does not claim to have excluded every conceivable coupling between rotation and gravity. It states that the evidence offered for one lacks the control that would make it evidence, and that the closest controlled tests came out null.

\begin{thebibliography}{99}

\bibitem{DP1974} B. E. DePalma, \emph{The Generation of a Unidirectional Force}, Simularity Institute, 1974.

\bibitem{DP1990} B. DePalma, On the Possibility of Extraction of Electrical Energy Directly from Space, \emph{Speculations in Science and Technology} 13 (1990) 283.

\bibitem{HT1989} H. Hayasaka and S. Takeuchi, Anomalous Weight Reduction on a Gyroscope's Right Rotations Around the Vertical Axis on the Earth, \emph{Physical Review Letters} 63 (1989) 2701--2704.

\bibitem{F1990} J. E. Faller, W. J. Hollander, P. G. Nelson and M. P. McHugh, Gyroscope-Weighing Experiment with a Null Result, \emph{Physical Review Letters} 64 (1990) 825--826.

\bibitem{QP1990} T. J. Quinn and A. Picard, The Mass of Spinning Rotors: No Dependence on Speed or Sense of Rotation, \emph{Nature} 343 (1990) 732--735.

\bibitem{T2022} P. Touboul et al., MICROSCOPE Mission: Final Results of the Test of the Equivalence Principle, \emph{Physical Review Letters} 129 (2022) 121102.

\bibitem{K1986} R. Kincheloe, Homopolar ``Free-Energy'' Generator Test, report on the Sunburst N machine, mid-1980s, as circulated in the DePalma literature, and discussed in \emph{The Skeptical Inquirer} 15 (1990--1991) 72.

\end{thebibliography}

\end{document}
